% ============================================================
% UNIDAD 2: AERONAVEGABILIDAD, CERTIFICACIÓN Y DOCUMENTACIÓN
% ============================================================
\clearpage
\section*{UNIDAD 2: AERONAVEGABILIDAD Y CERTIFICACIÓN}
\subsection*{Aeronaves, mantenimiento, partes y documentación técnica}

\begin{recuadro}
\textbf{Objetivo de la unidad:} comprender qué significa que una aeronave sea
\textbf{aeronavegable}, cómo se certifican el diseño, cada avión, el mantenimiento y las partes,
qué documentación respalda todo ese proceso y qué estándares (ATA, ISO, AN, MS, NAS, MIL) usa a
diario un técnico aviónico.
\end{recuadro}

% ------------------------------------------------------------
\section*{1. ¿QUÉ ES LA AERONAVEGABILIDAD?}

Una aeronave es \textbf{aeronavegable} cuando cumple dos condiciones a la vez:
\begin{enumerate}
  \item \textbf{Está conforme a su diseño aprobado} (certificado tipo), incluidas todas las
        modificaciones aprobadas; y
  \item \textbf{está en condiciones de operar en forma segura}: su desgaste, su estado y su
        mantenimiento son los adecuados.
\end{enumerate}
Si falta cualquiera de las dos, la aeronave no es aeronavegable, aunque ``vuele''. Un avión con un
transpondedor instalado sin aprobación, o con una inspección vencida, no es aeronavegable.

\paragraph*{La certificación es todo lo que hace que un avión, una persona o un taller estén
habilitados oficialmente para volar o para trabajar en aviación.} Conviene pensarla como
\textbf{tres cosas que se certifican} (la aeronave, las personas y las organizaciones) más la
\textbf{documentación} que lo respalda. Esta unidad trata la aeronave y la documentación; la
Unidad 3, las personas; la Unidad 4, las organizaciones.

% ------------------------------------------------------------
\section*{2. EL CICLO DE VIDA DE LA CERTIFICACIÓN DE UNA AERONAVE}

\begin{center}
\begin{tikzpicture}[node distance=0.35cm, every node/.style={font=\small, align=center},
  paso/.style={draw=ppAzul, very thick, rounded corners, fill=ppAzul!8, text width=2.6cm,
  minimum height=1.25cm}]
  \node[paso] (a) {\textbf{Diseño}\\ Certificado tipo};
  \node[paso, right=of a] (b) {\textbf{Producción}\\ Certificado de producción};
  \node[paso, right=of b] (c) {\textbf{Cada avión}\\ Certificado de aeronavegabilidad};
  \node[paso, right=of c] (d) {\textbf{Toda su vida}\\ Aeronavegabilidad continuada};
  \draw[-{Stealth}, thick] (a) -- (b);
  \draw[-{Stealth}, thick] (b) -- (c);
  \draw[-{Stealth}, thick] (c) -- (d);
\end{tikzpicture}
\end{center}

\begin{center}
\small
\begin{tabular}{|p{0.27\linewidth}|p{0.18\linewidth}|p{0.45\linewidth}|}
\hline
\textbf{Documento} & \textbf{Quién lo obtiene} & \textbf{Qué certifica} \\ \hline
\textbf{Certificado tipo} (con su hoja de datos, \textbf{TCDS}) & El fabricante & Que el
\textbf{diseño} de un modelo de avión, motor o hélice cumple las normas de aeronavegabilidad
(RAAC 21 y 23/25/27/29). \\ \hline
\textbf{Certificado tipo suplementario} (STC) & Quien diseña una modificación & Que una
\textbf{modificación mayor} al diseño (por ejemplo, instalar una nueva aviónica integrada) es
aceptable. \\ \hline
\textbf{Certificado de producción} & El fabricante & Que puede fabricar en serie aeronaves o
partes iguales al diseño aprobado. \\ \hline
\textbf{Certificado de aeronavegabilidad} & Cada avión individual & Que \textbf{ese avión}
coincide con el diseño aprobado y está en condiciones de volar seguro. Sigue válido solo si se
mantiene según el programa de mantenimiento aprobado. Puede ser estándar o especial (restringido,
experimental, deportivo liviano, permiso especial de vuelo). \\ \hline
\textbf{Directiva de aeronavegabilidad} (AD) & La emite la autoridad & Una inspección o
modificación \textbf{obligatoria} ante un problema descubierto en servicio. \\ \hline
\textbf{Aprobación para el retorno al servicio} & La firma una persona u organización habilitada
& Que, después de un trabajo de mantenimiento, el avión quedó aeronavegable. Esa firma tiene
\textbf{responsabilidad legal}. \\ \hline
\end{tabular}
\end{center}

\subsubsection*{2.1. Aeronavegabilidad continuada}
Una vez en servicio, la aeronave debe seguir siendo aeronavegable. Eso se logra con:
\begin{itemize}
  \item el \textbf{programa de mantenimiento} aprobado (inspecciones por horas de vuelo, ciclos o
        calendario);
  \item el cumplimiento de las \textbf{directivas de aeronavegabilidad (AD)};
  \item la evaluación de los \textbf{boletines de servicio (SB)} del fabricante: son
        recomendaciones, salvo que una AD los haga obligatorios;
  \item el control de los \textbf{componentes de vida limitada} (que deben retirarse al cumplir
        un número de horas o ciclos);
  \item los \textbf{registros} de todo lo realizado.
\end{itemize}

\begin{recuadro}
\textbf{Caso: UPS 2976 (Louisville, EE.\,UU., 4 de noviembre de 2025).} Un MD-11F de carga
perdió el motor izquierdo junto con su pilón durante el despegue y se estrelló; murieron los 3
tripulantes y 11 personas en tierra. La investigación de la NTSB encontró fisuras por fatiga en
un componente del montaje del pilón, y que el fabricante había documentado en 2011 fallas
anteriores de esa pieza en una carta de servicio, sin considerarlas un riesgo para la seguridad
del vuelo. Después del accidente, la FAA emitió \textbf{directivas de aeronavegabilidad de
emergencia} que obligaron a inspeccionar todos los MD-11 y DC-10 antes de volver a volar.
\textbf{Para discutir:} ¿qué diferencia hay entre una recomendación del fabricante y una AD? ¿Qué
rol juegan el operador y su programa de mantenimiento?
\end{recuadro}

% ------------------------------------------------------------
\section*{3. EL MANTENIMIENTO SEGÚN LA RAAC PARTE 43}

La \textbf{RAAC 43} define qué es mantenimiento y cómo debe hacerse y registrarse.

\begin{center}
\small
\begin{tabular}{|p{0.26\linewidth}|p{0.64\linewidth}|}
\hline
\textbf{Término} & \textbf{Significado} \\ \hline
Mantenimiento & Inspección, revisión general (\textit{overhaul}), reparación, conservación y
reemplazo de partes. \\ \hline
Mantenimiento preventivo & Tareas menores y simples de conservación y reemplazo de pequeñas
partes estándar, que no requieren operaciones de montaje complejas. La regulación las enumera.
\\ \hline
Reparación / alteración \textbf{mayor} & Puede afectar de manera apreciable el peso, el
balance, la resistencia estructural, el funcionamiento de los motores, las características de
vuelo u otras cualidades de aeronavegabilidad, o no puede hacerse con prácticas elementales.
Requiere \textbf{datos aprobados} y un formulario específico. \\ \hline
Reparación / alteración \textbf{menor} & Toda la que no es mayor. \\ \hline
Reconstrucción & Restaurar un producto a la condición de ``nuevo'', con partes nuevas o
equivalentes. \\ \hline
Retorno al servicio & Declaración, firmada por persona habilitada, de que el producto quedó
aeronavegable. \\ \hline
\end{tabular}
\end{center}

\paragraph*{¿Quién puede hacer mantenimiento y aprobar el retorno al servicio?}
Según los alcances de cada uno: el titular de una \textbf{licencia} de mecánico con las
habilitaciones correspondientes (RAAC 65), un \textbf{Taller Aeronáutico de Reparación} (RAAC 145)
a través de su representante técnico, un \textbf{explotador} con su propia organización de
mantenimiento (RAAC 121 o 135) y el \textbf{fabricante} para sus productos.

\paragraph*{El registro del trabajo.}
Todo trabajo debe quedar registrado en los historiales de la aeronave, motor, hélice o componente,
con al menos: \textbf{descripción del trabajo} (o referencia a datos aceptables), \textbf{fecha},
\textbf{nombre, firma y número de licencia} de quien aprueba el retorno al servicio. Las
reparaciones y alteraciones mayores se documentan además en un formulario específico (equivalente
al formulario 337 de la FAA). \textbf{Si no está registrado, para la ley no se hizo.}

% ------------------------------------------------------------
\section*{4. LOS DOCUMENTOS DE LA AERONAVE}

\subsubsection*{4.1. Documentos a bordo}
\begin{itemize}
  \item \textbf{Certificado de matrícula} (Registro Nacional de Aeronaves). La matrícula argentina
        empieza con \textbf{LV} (o LQ para aeronaves del Estado).
  \item \textbf{Certificado de aeronavegabilidad} vigente.
  \item \textbf{Manual de vuelo} aprobado (AFM) con sus limitaciones.
  \item \textbf{Licencia de la estación de radio} de a bordo.
  \item \textbf{Libro de a bordo} o registro técnico de vuelo, donde se anotan los vuelos y las
        fallas observadas.
  \item \textbf{Licencias y certificados médicos} de la tripulación; seguro obligatorio.
\end{itemize}

\subsubsection*{4.2. Historiales y registros de mantenimiento}
Historial de la célula, de cada motor y de cada hélice; registro del cumplimiento de AD;
control de componentes de vida limitada; tarjetas de trabajo (\textit{task cards}); peso y
balanceo actualizado; lista de equipos instalados.

\subsubsection*{4.3. Manuales técnicos}
\begin{center}
\small
\begin{tabular}{|p{0.16\linewidth}|p{0.74\linewidth}|}
\hline
\textbf{Sigla} & \textbf{Manual} \\ \hline
AMM & \textit{Aircraft Maintenance Manual}: procedimientos de mantenimiento de la aeronave. \\
\hline
IPC / IPB & \textit{Illustrated Parts Catalog}: catálogo ilustrado de partes con sus números de
parte (P/N) e intercambiabilidad. \\ \hline
WDM & \textit{Wiring Diagram Manual}: diagramas de cableado; el manual más usado por el técnico
aviónico. \\ \hline
SRM & \textit{Structural Repair Manual}: reparaciones estructurales. \\ \hline
CMM & \textit{Component Maintenance Manual}: mantenimiento de un componente en taller (por
ejemplo, un transceptor VHF). \\ \hline
MPD & \textit{Maintenance Planning Document}: tareas e intervalos para armar el programa de
mantenimiento. \\ \hline
MEL / MMEL & Lista de equipo mínimo: qué equipos pueden estar inoperativos y en qué condiciones
se puede despachar el avión igual. \\ \hline
FIM / TSM & Manual de aislamiento de fallas: guía de diagnóstico a partir de un mensaje o
síntoma. \\ \hline
\end{tabular}
\end{center}
Los manuales deben estar en su \textbf{última revisión}: trabajar con una copia vieja o
impresa ``para tenerla a mano'' es una causa clásica de errores.

% ------------------------------------------------------------
\section*{5. LAS PARTES AERONÁUTICAS}

Una parte solo puede instalarse en una aeronave si es una \textbf{parte aprobada} y está
acompañada de su \textbf{documentación de aeronavegabilidad}, que garantiza su origen y su estado:
\begin{itemize}
  \item \textbf{Formulario FAA 8130-3} o \textbf{EASA Form 1}: certificado de liberación
        autorizada de una parte nueva o reparada.
  \item \textbf{Certificado de conformidad} del fabricante.
  \item Partes fabricadas bajo \textbf{TSO} (orden técnica estándar), \textbf{PMA} (aprobación de
        fabricación de partes) o \textbf{partes estándar} reconocidas (AN, MS, NAS).
\end{itemize}

\paragraph*{Estados de una parte.} \textit{Serviciable} (lista para instalar, con su etiqueta y
certificado), \textit{inservible} o \textit{reparable} (con etiqueta roja o amarilla, separada
físicamente), \textit{chatarra} (debe inutilizarse para que no vuelva al circuito).

\begin{recuadro}
\textbf{Partes no aprobadas sospechosas (SUP).} Son partes de origen dudoso: falsificadas,
recuperadas de aviones accidentados sin control, con documentación adulterada o fuera de vida
útil. Están entre los riesgos más graves del mantenimiento. Señales de alerta: precio muy bajo,
documentación incompleta o fotocopiada, marcas borradas, vendedor desconocido. Ante la duda,
\textbf{la parte no se instala} y se reporta a la autoridad.
\end{recuadro}

% ------------------------------------------------------------
\section*{6. LOS ESTÁNDARES DE LA DOCUMENTACIÓN}

\subsubsection*{6.1. ATA 100 / ATA iSpec 2200}
La \textbf{especificación ATA 100} (hoy integrada en la \textbf{ATA iSpec 2200}) organiza toda la
documentación técnica en \textbf{capítulos numerados} iguales para todos los fabricantes. Así, el
sistema de comunicaciones está en el capítulo 23 tanto en un Boeing como en un Airbus o un
Embraer. La referencia completa tiene tres pares de números: \textbf{capítulo -- sección -- tema}.
Por ejemplo, \texttt{23-11-00} es el sistema de HF dentro de Comunicaciones.

\begin{center}
\small
\begin{tabular}{|c|p{0.32\linewidth}||c|p{0.32\linewidth}|}
\hline
\textbf{Cap.} & \textbf{Sistema} & \textbf{Cap.} & \textbf{Sistema} \\ \hline
05 & Límites de tiempo, inspecciones & 31 & \textbf{Instrumentos indicadores y registradores} \\
\hline
12 & Servicio (carga de fluidos) & 32 & Tren de aterrizaje \\ \hline
20 & Prácticas estándar & 33 & \textbf{Luces} \\ \hline
21 & Aire acondicionado & 34 & \textbf{Navegación} \\ \hline
\textbf{22} & \textbf{Piloto automático} & 42 & \textbf{Aviónica modular integrada} \\ \hline
\textbf{23} & \textbf{Comunicaciones} & 44 & \textbf{Sistemas de cabina (entretenimiento)} \\
\hline
\textbf{24} & \textbf{Energía eléctrica} & 45 & \textbf{Sistema central de mantenimiento} \\ \hline
26 & Protección contra incendio & 46 & \textbf{Sistemas de información} \\ \hline
28 & Combustible & 49 & APU (unidad de potencia auxiliar) \\ \hline
30 & Protección contra hielo y lluvia & 71--80 & Motor \\ \hline
\end{tabular}
\end{center}

\subsubsection*{6.2. ATA 104}
La especificación \textbf{ATA 104} define cómo se estructura la \textbf{capacitación} del personal
de mantenimiento en niveles crecientes: desde la familiarización general con la aeronave hasta
el nivel de mantenimiento de línea y base, y la especialización en sistemas o componentes. Es la
base de los cursos de tipo que dan los fabricantes y los centros de instrucción.

\subsubsection*{6.3. Normas de partes y componentes: AN, MS, NAS, MIL}
\begin{center}
\small
\begin{tabular}{|p{0.13\linewidth}|p{0.36\linewidth}|p{0.40\linewidth}|}
\hline
\textbf{Norma} & \textbf{Origen} & \textbf{Ejemplos} \\ \hline
\textbf{AN} & \textit{Army--Navy}: normas militares conjuntas de EE.\,UU. (desde los años 40). &
Bulones AN3 a AN20, tuercas AN365, arandelas AN960. \\ \hline
\textbf{MS} & \textit{Military Standard}: reemplazó a muchas normas AN. & Remaches MS20470,
tornillos MS24693, terminales y empalmes de cableado. \\ \hline
\textbf{NAS} & \textit{National Aerospace Standard}: normas de la industria aeroespacial de EE.\,UU.
& Tornillos NAS1352, bulones de alta resistencia NAS6203. \\ \hline
\textbf{MIL} & \textit{Military Specifications} y estándares militares. & Cable MIL-W-22759,
conectores MIL-DTL-38999, bus de datos \textbf{MIL-STD-1553}. \\ \hline
\end{tabular}
\end{center}
Las normas fijan material, dimensiones, tolerancias, tratamientos y ensayos. Por eso un tornillo
aeronáutico \textbf{nunca se reemplaza ``por uno parecido''}, sino por el \textbf{número de parte
exacto} del catálogo ilustrado.

\begin{recuadro}
\textbf{Caso: British Airways 5390 (Reino Unido, 1990).} Al reemplazar el parabrisas del piloto
de un BAC 1-11, un supervisor eligió los tornillos ``a ojo'' en el depósito, sin consultar el
catálogo: 84 de los 90 tornillos eran apenas más finos que los especificados. En vuelo, la presión
de cabina arrancó el parabrisas y el comandante quedó parcialmente fuera de la cabina; la
tripulación logró aterrizar y no hubo víctimas fatales. El caso muestra el valor del número de
parte, del catálogo y de la verificación independiente.
\end{recuadro}

\subsubsection*{6.4. Otros estándares que usa un técnico aviónico}
\begin{itemize}
  \item \textbf{ARINC:} estándares de la industria para aviónica (por ejemplo, el bus de datos
        \textbf{ARINC 429} y los formatos de equipos y bandejas).
  \item \textbf{RTCA DO-160:} ensayos ambientales de equipos de a bordo (temperatura, vibración,
        interferencia electromagnética). \textbf{DO-178C} y \textbf{DO-254}: desarrollo de
        software y hardware de aviónica.
  \item \textbf{ISO 9001} (gestión de la calidad) y sus versiones aeroespaciales \textbf{AS/EN
        9100} (fabricantes), \textbf{AS9110} (organizaciones de mantenimiento) y \textbf{AS9120}
        (distribuidores). Se conectan con el sistema de calidad que exige la RAAC 145.
\end{itemize}

% ------------------------------------------------------------
\section*{7. CÓMO SE CONECTA TODO}
Cada trabajo de mantenimiento recorre la misma cadena:
\begin{enumerate}
  \item Un técnico \textbf{con licencia (RAAC 65)}
  \item trabaja en un \textbf{taller certificado (RAAC 145)},
  \item sigue el \textbf{manual del fabricante, organizado por capítulos ATA},
  \item usa \textbf{partes aprobadas, con su documentación} y que cumplen las normas AN, MS, NAS
        o MIL,
  \item registra el trabajo y \textbf{firma el retorno al servicio} (RAAC 43),
  \item y así el avión mantiene su \textbf{certificado de aeronavegabilidad}.
\end{enumerate}

% ------------------------------------------------------------
\section*{8. ACTIVIDADES DE LA UNIDAD}
\begin{enumerate}
  \item \textbf{Recorrido de certificación:} tomar un trabajo de aviónica (por ejemplo, reemplazar
        un receptor VOR) y seguir todo el recorrido: ubicar el capítulo ATA del sistema;
        verificar el número de parte del componente en el catálogo; revisar si hay directivas
        de aeronavegabilidad aplicables; completar el registro de mantenimiento; indicar quién
        puede firmar el retorno al servicio y con qué licencia.
  \item \textbf{Verificación de aeronavegabilidad simulada:} con un juego de documentos de una
        aeronave (reales o preparados por el docente), verificar documentos a bordo, vencimientos
        de inspecciones, cumplimiento de AD y componentes de vida limitada. Informar si la
        aeronave es aeronavegable y por qué.
  \item \textbf{¿Mayor o menor?} Clasificar una lista de trabajos (cambiar una lámpara, instalar un
        GPS con pantalla nueva en el panel, reparar un cable, reemplazar un transpondedor por otro
        del mismo P/N) en mantenimiento preventivo, alteración menor o mayor.
  \item \textbf{Tabla comparativa de estándares} ISO, AN, MS, NAS y MIL, con su aplicación en
        mantenimiento y documentación técnica, y un ejemplo de pieza del taller de la escuela.
  \item \textbf{Partes sospechosas:} con tres certificados de partes (uno correcto, uno incompleto,
        uno adulterado), decidir cuál se acepta y redactar el informe de rechazo.
\end{enumerate}
